\documentclass[12pt, legalpaper]{article}

\usepackage[margin=3cm]{geometry}
\usepackage{natbib}
\usepackage{lipsum}
\usepackage{algorithm}
\usepackage{algorithmic}
\usepackage{amsfonts}
\usepackage{graphicx}
\usepackage{epstopdf}
\usepackage{algorithm}
\bibliographystyle{unsrt} % or another style like unsrt, abbrv, etc.

\usepackage{xcolor}

\usepackage{color}

\newcommand{\rr}[1]{\textcolor{red}{#1}}
\newcommand{\rrr}[2]{\textcolor{blue}{(#1)}\textnormal{$\leftarrow$}\textcolor{red}{#2}}
%\newcommand{\rrrr}[1]{\textcolor{blue}{(#1) $\leftarrow$ (remove)}}
\usepackage[normalem]{ulem}\newcommand{\rrrr}[1]{\textcolor{blue}{\sout{(#1)}}}
% \newcommand{\rrrr}[1]{\textcolor{blue}{(#1)}\textnormal{$\leftarrow$}\textcolor{blue}{{\rm remove}}}
\newcommand{\bb}[1]{\textcolor{blue}{#1}}
\newcommand{\bbb}[2]{\textcolor{blue}{(#1) $\leftarrow$ #2}}
\newcommand{\bbbb}[1]{\textcolor{blue}{(#1) $\leftarrow$ (remove)}}
\newcommand{\ignore}[1]{}

%\usepackage{algorithm}\usepackage{algorithmic}

\usepackage{amsmath}
\usepackage{amssymb}

\usepackage{graphicx} % Required for inserting images
\usepackage{hyperref}
\usepackage{cleveref}

\usepackage{datetime}

\newdateformat{monthyeardate}{%
  \monthname[\THEMONTH], \THEYEAR}




\title{TIG CUR Decomposition Challenge}
\author{TIG Labs - Yuji Nakatsukasa}
\date{\monthyeardate\today}

\begin{document}

\maketitle

% \begin{abstract}
%    We propose an ...
% \end{abstract}


%%%
\section{Introduction}



In the era of big data, analysts routinely work with matrices containing tens or even hundreds of thousands of rows and columns. This scale is computationally taxing; however, the data can contain redundancy. Isolating the informative low-dimensional structure has therefore become a critical first step in any analysis. Matrix factorizations, in particular, low-rank decompositions and low-rank approximations sit at the heart of modern data science and numerical analysis, providing the theoretical backbone for fast solvers in scientific computing. Within this landscape, CUR approximations have gained popularity because they preserve the actual rows and columns of the original matrix, yielding interpretable, sparsity-aware low rank approximations. CUR approximations have been used in a variety of areas including genetic data \cite{bodor2012rcur}, imaging \cite{abdelgawad2024efficient}, and recommender systems \cite{mahoney2006tensor}.

\section{The TIG Challenge}
A challenge instance is defined by the track parameters
\[
m,\qquad n,\qquad \Delta,
\]
a singular-value structure, and a random seed. From these parameters, the
generator constructs $S=8$ related CUR sub-instances as described in
Section~\ref{sec: instancegen}. The challenge proceeds in three steps:
\begin{itemize}
    \item \textbf{Step 1. Generation.} The generator computes shared
    orthonormal bases $U$ and $V$ once, then uses them to construct
    \[
    (A_1,k_1),\ldots,(A_S,k_S),
    \]
    where $A_s\in\mathbb{R}^{m\times n}$ has true rank $r_s$ and $k_s$ is
    the target rank of the CUR approximation. The true ranks and
    target-rank ratios are sampled using the stratified procedure described
    in Section~\ref{sec: instancegen}.
    \item \textbf{Step 2. Solving.} For each sub-instance $(A_s,k_s)$, the
    innovator's algorithm selects exactly $k_s$ columns and $k_s$ rows of
    $A_s$. Let $\mathcal F$ contain the four sub-instances with the largest
    target ranks, with ties resolved in favour of the lower sub-instance
    index. For
    $s\in\mathcal F$, the innovator uses the challenge's fixed fast linking
    matrix calculation when searching and assessing solution quality, but
    returns only the row and column indices. For $s\notin\mathcal F$, the
    innovator also returns a finite linking matrix $U_s^{\mathrm{link}}$.
    \item \textbf{Step 3. Verification and scoring.} For
    $s\in\mathcal F$, the verifier reconstructs $U_s^{\mathrm{link}}$ using
    exactly the same fixed fast calculation. It uses the submitted linking
    matrix for the other four sub-instances. Each resulting approximation
    $CUR_s=C_sU_s^{\mathrm{link}}R_s$ is scored by the continuous rule in
    Section~\ref{sec: score}, and the final score is the arithmetic mean over
    all eight sub-instances.
\end{itemize}
The use of multiple sub-instances allows the expensive instance-generation
work to be amortised while testing an innovator's algorithm across a range
of true-rank and target-rank regimes.


%%

%%%# Sub-instance Generation

\section{Sub-instance Generation}\label{sec: instancegen}

An instance is specified by a random seed, matrix dimensions $m,n$, a
coherence parameter $\Delta$, and a choice of singular-value structure.
Let
\[
\tau=\min\{m,n\}.
\]

Using the random seed, generate Gaussian matrices
\[
G_U\in\mathbb{R}^{m\times\tau},
\qquad
G_V\in\mathbb{R}^{n\times\tau}.
\]
The columns are scaled
using the diagonal matrix $D$ determined by $\Delta$, and QR decompositions
are computed:
\[
G_UD=UR_U,
\qquad
G_VD=VR_V.
\]
This gives orthonormal bases
\[
U=(u_1,\ldots,u_\tau),
\qquad
V=(v_1,\ldots,v_\tau).
\]
The costly QR decompositions are therefore performed only once and the
resulting pools of singular vectors are reused to construct all
sub-instances.\\

We generate $S=8$ sub-instances.

\subsection{True ranks}

For each sub-instance $s$, define
\[
\alpha_s=\frac{r_s}{\tau},
\]
where $r_s$ is the true rank of the matrix. We sample one $\alpha_s$ from
each of the stratified intervals
\[
\mathcal A_1=[0.03,0.055],\;
\mathcal A_2=[0.055,0.08],\;
\mathcal A_3=[0.08,0.105],\;
\mathcal A_4=[0.105,0.13],
\]
\[
\mathcal A_5=[0.13,0.155],\;
\mathcal A_6=[0.155,0.18],\;
\mathcal A_7=[0.18,0.20],\;
\mathcal A_8=[0.20,0.22].
\]
and set
\[
r_s=\left\lfloor\alpha_s\tau\right\rceil.
\]

Thus every challenge instance contains sub-instances spanning a range of
true ranks, while the precise ranks vary with the random seed.

\subsection{Randomly overlapping singular-vector sets}

For each sub-instance independently sample
\[
I_s\subseteq\{1,\ldots,\tau\},
\qquad |I_s|=r_s,
\]
uniformly without replacement.

The sets $I_s$ are not required to be disjoint. Consequently different
sub-instances naturally share some of the singular directions generated by
the common QR factorisations.

For two sub-instances $s,t$, by the law of total expectation,
\[
\mathbb E\!\left[|I_s\cap I_t|\right]
=
\frac{r_sr_t}{\tau}.
\]
Equivalently, the expected overlap as a fraction of the smaller
sub-instance is
\[
\frac{\mathbb E[|I_s\cap I_t|]}
     {\min\{r_s,r_t\}}
=
\max\{\alpha_s,\alpha_t\}.
\]
The chosen $\alpha$ ranges therefore control the typical amount of
structural information shared between sub-instances. So over all sub instances the maximum expected pairwise overlap as a fraction of the smaller of the two instances is given by $\max_{s\in \{1,\ldots,S\}} \alpha_s$.\footnote{We are using $\frac{\mathbb E[|I_s\cap I_t|]}
     {\min\{r_s,r_t\}}$ as a measure of how much information is shared between subinstances. There may be better measures.}

\subsection{Singular values and random re-pairing}

For each sub-instance, write
\[
I_s=\{i_1,\ldots,i_{r_s}\}.
\]

First generate a base singular-value profile
\[
\bar{\sigma}^{(s)}_j
=
f_{\mathrm{type}}(r_s,j),
\qquad j=1,\ldots,r_s,
\]
using the singular-value structure specified by the instance.

To avoid all sub-instances with the same singular-value type having
identical spectra up to rank, independently perturb these values:
\[
\sigma^{(s)}_j
=
\bar{\sigma}^{(s)}_j
\left(1+\varepsilon^{(s)}_j\right),
\qquad
\varepsilon^{(s)}_j
\sim
\operatorname{Uniform}(-0.15,0.15).
\]
Thus each singular value is independently perturbed by up to $15\%$,
giving each sub-instance a slightly different spectrum while preserving
the overall singular-value structure.

Next generate an independent random permutation
\[
\pi_s:\{1,\ldots,r_s\}\rightarrow\{1,\ldots,r_s\}.
\]
The selected left and right singular vectors are then randomly re-paired
when constructing the matrix:
\[
A_s
=
\sum_{j=1}^{r_s}
\sigma^{(s)}_j
u_{i_j}
v_{i_{\pi_s(j)}}^T.
\]

Equivalently,
\[
A_s
=
U[:,I_s]\,
\Sigma_s\,
V[:,I_s^{\pi_s}]^T,
\]
where
\[
\Sigma_s
=
\operatorname{diag}
\left(
\sigma^{(s)}_1,\ldots,\sigma^{(s)}_{r_s}
\right).
\]

Since both selected bases remain orthonormal,
\[
\operatorname{rank}(A_s)=r_s.
\]

Thus sub-instances may share some underlying singular vectors, allowing the
cost of the original QR factorisations to be amortised, while the random
index sets, perturbed singular values, and independently randomised
left--right pairings make the resulting matrices structurally different.

\subsection{Target ranks}

For each sub-instance define
\[
\rho_s=\frac{k_s}{r_s},
\]
where $k_s$ is the target rank of the CUR approximation.

Sample one value from each of
\[
\mathcal R_1=[0.10,0.19],\;
\mathcal R_2=[0.19,0.28],\;
\mathcal R_3=[0.28,0.37],\;
\mathcal R_4=[0.37,0.46],
\]
\[
\mathcal R_5=[0.46,0.55],\;
\mathcal R_6=[0.55,0.64],\;
\mathcal R_7=[0.64,0.72],\;
\mathcal R_8=[0.72,0.80].
\]

The eight sampled values are randomly permuted before being paired with the
eight true ranks, preventing a fixed relationship between $\alpha_s$ and
$\rho_s$. Finally, set
\[
k_s=\left\lfloor\rho_s r_s\right\rceil.
\]

The resulting challenge instance consists of
\[
(A_1,k_1),\ldots,(A_S,k_S).
\]


\section{Linking Matrices and Solution Size}\label{sec:linking-matrices}

Once the column and row indices have been selected, the CUR approximation
still requires a linking matrix $U_s^{\mathrm{link}}\in
\mathbb{R}^{k_s\times k_s}$.  There are two natural ways to compute this
matrix.  A fast method uses thin QR factorisations
\[
C_s=Q_C T_C,
\qquad
R_s^T=Q_R T_R,
\]
followed by triangular solves to form
\[
U_{s,\mathrm{fast}}^{\mathrm{link}}
=
T_C^{-1}\bigl(Q_C^T A_s Q_R\bigr)T_R^{-T}.
\]
A slower, more accurate method uses rank-revealing SVD pseudoinverses to
approximate the optimal linking matrix
\[
U_s^*=C_s^+A_sR_s^+.
\]
The fast method is substantially cheaper, whereas the SVD method is more
robust when the selected rows or columns are poorly conditioned.

Returning every linking matrix can consume a large part of the TIG solution
limit.  In single-precision format, one linking matrix requires
\[
4k_s^2\ \text{bytes},
\]
before serialisation overhead. To reduce this output size, let $\mathcal F$
contain the four sub-instances having the largest target ranks $k_s$, with
ties resolved in favour of the lower sub-instance index. These are the four
largest linking matrices because their storage grows as $k_s^2$.

For $s\in\mathcal F$, the innovator must use the fixed fast QR calculation
above when evaluating candidate row and column sets and when calculating the
quality of its final selection. It returns only the selected column and row
index sets $(\mathcal J_s,\mathcal I_s)$; it does not return the resulting
linking matrix. The verifier reconstructs
\[
C_s=A_s[:,\mathcal J_s],
\qquad
R_s=A_s[\mathcal I_s,:],
\]
and computes $U_{s,\mathrm{fast}}^{\mathrm{link}}$ with exactly the same
routine. The challenge implementation provides this routine to both the
innovator and verifier, including its arithmetic precision, QR convention,
and solve order. Thus the innovator can calculate the same CUR approximation
and solution quality that the verifier will score; the verifier's result
remains authoritative.

For each of the remaining four sub-instances, $s\notin\mathcal F$, the
innovator returns $U_s^{\mathrm{link}}$ as well as the two index sets and may
use the slower, more accurate SVD construction. Choosing the four largest
$k_s$ values for verifier-side reconstruction maximises the reduction in
returned data while leaving linking-matrix construction under innovator
control for the other four sub-instances.





\section{Scoring}\label{sec: score}

Each challenge instance consists of multiple sub-instances, where sub-instance $i$
is given by a matrix $A_i$ and target rank $k_i$. For each sub-instance, define
the approximation-quality ratio
\begin{equation}\label{eq err}
q_i
:=
\frac{\|A_i-CUR_i\|_F}
     {\|A_i-SVD_{k_i}(A_i)\|_F},
\end{equation}
where $SVD_{k_i}(A_i)$ denotes the optimal rank-$k_i$ approximation to $A_i$.

Since the truncated SVD is the optimal rank-$k_i$ approximation, $q_i=1$
corresponds to SVD-quality performance, while larger values of $q_i$ indicate
worse approximations. From the CUR approximation literature, the value
\[
k_i+1
\]
provides a natural failure threshold.

Because the target ranks $k_i$ vary between sub-instances, the raw values
$q_i$ are not directly on the same scale. In particular, the failure threshold
itself changes with $k_i$. We therefore normalise the approximation quality by
defining\footnote{We use a max below in case numerical error decreases $q_i$ below 1.}
\begin{equation}\label{eq normalized_quality}
z_i
:=
\frac{\log(\max\{q_i,1\})}
     {\log(k_i+1)}.
\end{equation}

Thus $z_i$ measures the multiplicative degradation from SVD-quality relative
to the literature threshold, on a scale that is comparable across different
target ranks.

The score for sub-instance $i$ is then defined as
\begin{equation}\label{eq scoring}
s_i
:=
\frac{1}{1+z_i}.
\end{equation}

Hence
\[
q_i=1
\quad\Longrightarrow\quad
s_i=1,
\]
and
\[
q_i=k_i+1
\quad\Longrightarrow\quad
s_i=\frac{1}{2}.
\]
Results better than the literature threshold therefore receive scores in
$(1/2,1]$, while results worse than the threshold receive scores below $1/2$.
The scoring rule continuously differentiates between solutions without
clipping all successful or unsuccessful solutions to the same value.

The final score for a challenge instance is the average over its $S$
sub-instances:
\[
\mathrm{Score}
=
\frac{1}{S}
\sum_{i=1}^{S}s_i.
\]



%%%
\section{Challenge Tracks}\label{sec:diffparams}
Within each challenge, there are various challenges tracks. These can range over instance size and/or type. Currently, the challenge supports varying sizes of matrices as well as the singular value structure.
\begin{itemize}

    \item \texttt{Track \{ m: 2000, n: 3000, poly: false \}}
    \item \texttt{Track \{ m: 4000, n: 4000, poly: false \}}
    \item \texttt{Track \{ m: 2000, n: 2000, poly: true \}}
    \item \texttt{Track \{ m: 2000, n: 3000, poly: true \}}
    \item \texttt{Track \{ m: 8000, n: 8000, poly: true \}}
\end{itemize}


 The coherence $\Delta$ is fixed but could be varied in the future to adapt the difficulty. Moreover we could introduce sparsity too.

Here we give our choice for the decay functions $f$, for a constant $a\in \mathbb{R}_+$ :
\begin{itemize}
    \item Exponential decay
    \begin{equation}
   f(\tau,j,a)=\exp\Bigg(-\frac{a\sqrt{j+1}}{\sqrt{\tau+1}}\Bigg)
    \end{equation}

  Note the $\sqrt{\tau}$ to stop the majority of $\sigma$ being $0$ as $\tau$ increased.

\item
Polynomial decaying singular values
\begin{equation}
   f(\tau,j,a)=\frac{c}{\left(\frac{j}{\tau}\right)^{2}+c}
\end{equation}
for $c=-\frac{e^{-a}}{e^{-a}-1}$ to match the exponential decay.
\end{itemize}

\section{Asymmetry Characteristics}\label{sec:asymmetry} To maintain the asymmetry required by the TIG protocol, the total time required by an innovator to solve the sub-instances should be significantly greater than the time required to generate the instance and verify the submitted solution. The sub-instance construction is designed to improve this asymmetry by amortising the expensive generation of the orthonormal bases across several CUR problems. In particular, the Gaussian matrices and QR decompositions \[ G_UD=UR_U, \qquad G_VD=VR_V \] are computed only once. The resulting pools of singular vectors are then reused to construct all $S$ sub-instances. The singular-vector index sets $I_s$ are sampled independently and are allowed to overlap. This is important for generation efficiency: the total number of singular directions used across all sub-instances may exceed $\tau$ without requiring a larger QR factorisation. At the same time, each sub-instance uses a different random subset of these directions, independently perturbed singular values, and an independently randomised pairing between its left and right singular vectors. Consequently, some structural information is necessarily shared between sub-instances. Verification is also simplified by the fact that the singular values of every generated matrix are known. Let \[ \sigma^{(s)}_{(1)} \geq \sigma^{(s)}_{(2)} \geq \cdots \geq \sigma^{(s)}_{(r_s)} \] denote the perturbed singular values of $A_s$, ordered from largest to smallest. For target rank $k_s$, the optimal rank-$k_s$ approximation error is therefore \[ \|A_s-SVD_{k_s}(A_s)\|_F = \left( \sum_{j=k_s+1}^{r_s} \left(\sigma^{(s)}_{(j)}\right)^2 \right)^{1/2}. \] Hence the denominator of the scoring rule can be calculated directly from the instance-generation data and does not require computing an SVD during verification.

\section{Pseudo Code}
This section contains pseudocode for the three main components of the
challenge:
\[
\textbf{generate} \;\longrightarrow\; \textbf{TIG harness}
\;\longrightarrow\; \textbf{verify}.
\]

\begin{algorithm}
\caption{\textbf{Generate} the eight CUR sub-instances.}
\label{sec:algorithm_generate}
\begin{algorithmic}[1]
\STATE \textbf{Input:} $m,n,\Delta$, singular-value type $f$, spectral
perturbation $\varepsilon_\sigma$, and random seed
\STATE \textbf{Output:} $S=8$ matrix--target-rank pairs $(A_s,k_s)$
\STATE Initialise all randomness from the supplied seed
\STATE $\tau \gets \min\{m,n\}$
\STATE Sample Gaussian matrices $G_U\in\mathbb{R}^{m\times\tau}$ and
$G_V\in\mathbb{R}^{n\times\tau}$
\STATE Construct $D=\mathrm{diag}(d_1,\ldots,d_\tau)$ with entries decreasing
linearly from $\Delta$ to $1$
\STATE Compute $G_UD=UR_U$ and $G_VD=VR_V$
\STATE Define $\mathcal A_1,\ldots,\mathcal A_8$ and
$\mathcal R_1,\ldots,\mathcal R_8$ as in Section~\ref{sec: instancegen}
\FOR{$s=1,\ldots,S$}
    \STATE Sample $\alpha_s\sim\mathrm{Uniform}(\mathcal A_s)$ and set
    $r_s\gets\lfloor\alpha_s\tau\rceil$
    \STATE Sample $\widetilde{\rho}_s\sim\mathrm{Uniform}(\mathcal R_s)$
\ENDFOR
\STATE Randomly permute $(\widetilde{\rho}_1,\ldots,\widetilde{\rho}_S)$ to
obtain $(\rho_1,\ldots,\rho_S)$
\STATE $\mathrm{SubInstances}\gets ()$
\FOR{$s=1,\ldots,S$}
    \STATE Sample $I_s\subseteq\{1,\ldots,\tau\}$ uniformly without
    replacement, with $|I_s|=r_s$
    \FOR{$j=1,\ldots,r_s$}
        \STATE $\bar{\sigma}^{(s)}_j \gets f(r_s,j)$
        \STATE Sample $\varepsilon^{(s)}_j \sim
        \mathrm{Uniform}(-\varepsilon_\sigma,\varepsilon_\sigma)$
        \STATE $\sigma^{(s)}_j \gets \bar{\sigma}^{(s)}_j
        (1+\varepsilon^{(s)}_j)$
    \ENDFOR
    \STATE Sample an independent random permutation $\pi_s$ of
    $\{1,\ldots,r_s\}$
    \STATE Write $I_s=\{i_1,\ldots,i_{r_s}\}$
    \STATE Construct
    $A_s \gets \sum_{j=1}^{r_s}\sigma^{(s)}_j
    u_{i_j}v_{i_{\pi_s(j)}}^T$
    \STATE $k_s\gets\lfloor\rho_s r_s\rceil$
    \STATE $\mathrm{SubInstances.append}\big((A_s,k_s)\big)$
\ENDFOR
\STATE \textbf{return} $\mathrm{SubInstances}$
\end{algorithmic}
\end{algorithm}

\begin{algorithm}
\caption{\textbf{TIG harness} that runs innovator submissions.}
\label{sec:algorithm_harness}
\begin{algorithmic}[1]
\STATE \textbf{Input:} $S=8$ sub-instances $(A_s,k_s)$
\STATE Let $\mathcal F$ be the four indices with the largest $k_s$, resolving
ties in favour of the lower sub-instance index
\STATE $\mathrm{Solution}\gets ()$
\FOR{$s=1,\ldots,S$}
    \STATE $(\mathcal J_s,\mathcal I_s)\gets
    \mathrm{Innovator\_Select}(A_s,k_s)$, using the prescribed linking rule
    for $s$
    \STATE $C_s\gets A_s[:,\mathcal J_s]$ and
    $R_s\gets A_s[\mathcal I_s,:]$
    \IF{$s\in\mathcal F$}
        \STATE $U_s^{\mathrm{link}}\gets
        \mathrm{FixedFastLink}(A_s,C_s,R_s)$ for the innovator's own quality
        calculation
        \STATE $\mathrm{Solution.append}\big((\mathcal J_s,\mathcal I_s)\big)$
    \ELSE
        \STATE $U_s^{\mathrm{link}}\gets
        \mathrm{Innovator\_Link}(A_s,C_s,R_s)$
        \STATE $\mathrm{Solution.append}
        \big((\mathcal J_s,\mathcal I_s,U_s^{\mathrm{link}})\big)$
    \ENDIF
\ENDFOR
\STATE \textbf{return} $\mathrm{Solution}$
\end{algorithmic}
\end{algorithm}

For every sub-instance, $\mathcal J_s$ contains the selected column indices
and $\mathcal I_s$ contains the selected row indices. For
$s\in\mathcal F$, the innovator must call the same
$\mathrm{FixedFastLink}$ implementation used by the verifier whenever it
calculates or reports the quality of its selection, but omits the resulting
matrix from its solution. For $s\notin\mathcal F$, it submits its chosen
$U_s^{\mathrm{link}}$. The matrices $C_s$ and $R_s$ are never returned
explicitly.

\begin{algorithm}
\caption{\textbf{Verification} that scores the submitted solutions.}
\label{sec:algorithm_verify}
\begin{algorithmic}[1]
\STATE \textbf{Input:} $S=8$ sub-instances $(A_s,k_s)$ and submitted
solutions
\STATE Let $\mathcal F$ be the four indices with the largest $k_s$, resolving
ties in favour of the lower sub-instance index
\STATE $\mathrm{TotalScore}\gets0$
\FOR{$s=1,\ldots,S$}
    \IF{$s\in\mathcal F$}
        \STATE $(\mathcal J_s,\mathcal I_s)\gets\mathrm{Solution}[s]$
    \ELSE
        \STATE $(\mathcal J_s,\mathcal I_s,U_s^{\mathrm{link}})
        \gets\mathrm{Solution}[s]$
        \STATE Require a finite $k_s\times k_s$ submitted linking matrix
    \ENDIF
    \STATE Require $k_s$ distinct in-bounds row indices and $k_s$ distinct
    in-bounds column indices
    \STATE $C_s\gets A_s[:,\mathcal J_s]$ and
    $R_s\gets A_s[\mathcal I_s,:]$
    \IF{$s\in\mathcal F$}
        \STATE $U_s^{\mathrm{link}}\gets
        \mathrm{FixedFastLink}(A_s,C_s,R_s)$
    \ENDIF
    \STATE $CUR_s\gets C_sU_s^{\mathrm{link}}R_s$
    \STATE $E_s^*\gets
    \left(\sum_{j=k_s+1}^{r_s}(\sigma_{(j)}^{(s)})^2\right)^{1/2}$
    \STATE $q_s\gets\|A_s-CUR_s\|_F/E_s^*$
    \STATE $z_s\gets\log(\max\{q_s,1\})/\log(k_s+1)$
    \STATE $\mathrm{TotalScore}\gets
    \mathrm{TotalScore}+1/(1+z_s)$
\ENDFOR
\STATE \textbf{return} $\mathrm{TotalScore}/S$
\end{algorithmic}
\end{algorithm}

\bibliography{references} % references.bib is the file containing the bib entries


\end{document}
